\documentclass[12pt,a4paper]{article}

\usepackage[T2A]{fontenc}
\usepackage[utf8]{inputenc}
\usepackage[russian]{babel}
\usepackage{amsmath, amssymb, amsthm}
\usepackage{geometry}
\usepackage{xcolor}
\usepackage{listings}
\usepackage{graphicx}
\usepackage{float}

\geometry{left=25mm, right=15mm, top=20mm, bottom=20mm}

\definecolor{codegreen}{rgb}{0,0.6,0}
\definecolor{codegray}{rgb}{0.5,0.5,0.5}
\definecolor{codepurple}{rgb}{0.58,0,0.82}
\definecolor{backcolour}{rgb}{0.98,0.98,0.98}

\lstdefinestyle{sqlstyle}{
    backgroundcolor=\color{backcolour},
    commentstyle=\color{codegreen},
    keywordstyle=\color{blue}\bfseries,
    numberstyle=\tiny\color{codegray},
    stringstyle=\color{codepurple},
    basicstyle=\ttfamily\footnotesize,
    breakatwhitespace=false,
    breaklines=true,
    captionpos=b,
    keepspaces=true,
    numbers=left,
    numbersep=5pt,
    showspaces=false,
    showstringspaces=false,
    showtabs=false,
    tabsize=2,
    language=SQL,
    literate=
      {А}{{\CYRA}}1
      {Б}{{\CYRB}}1
      {В}{{\CYRV}}1
      {Г}{{\CYRG}}1
      {Д}{{\CYRD}}1
      {Е}{{\CYRE}}1
      {Ё}{{\CYRYO}}1
      {Ж}{{\CYRZH}}1
      {З}{{\CYRZ}}1
      {И}{{\CYRI}}1
      {Й}{{\CYRISHRT}}1
      {К}{{\CYRK}}1
      {Л}{{\CYRL}}1
      {М}{{\CYRM}}1
      {Н}{{\CYRN}}1
      {О}{{\CYRO}}1
      {П}{{\CYRP}}1
      {Р}{{\CYRR}}1
      {С}{{\CYRS}}1
      {Т}{{\CYRT}}1
      {У}{{\CYRU}}1
      {Ф}{{\CYRF}}1
      {Х}{{\CYRH}}1
      {Ц}{{\CYRC}}1
      {Ч}{{\CYRCH}}1
      {Ш}{{\CYRSH}}1
      {Щ}{{\CYRSHCH}}1
      {Ъ}{{\CYRHRDSN}}1
      {Ы}{{\CYRERY}}1
      {Ь}{{\CYRSFTSN}}1
      {Э}{{\CYREREV}}1
      {Ю}{{\CYRYU}}1
      {Я}{{\CYRYA}}1
      {а}{{\cyra}}1
      {б}{{\cyrb}}1
      {в}{{\cyrv}}1
      {г}{{\cyrg}}1
      {д}{{\cyrd}}1
      {е}{{\cyre}}1
      {ё}{{\cyryo}}1
      {ж}{{\cyrzh}}1
      {з}{{\cyrz}}1
      {и}{{\cyri}}1
      {й}{{\cyrishrt}}1
      {к}{{\cyrk}}1
      {л}{{\cyrl}}1
      {м}{{\cyrm}}1
      {н}{{\cyrn}}1
      {о}{{\cyro}}1
      {п}{{\cyrp}}1
      {р}{{\cyrr}}1
      {с}{{\cyrs}}1
      {т}{{\cyrt}}1
      {у}{{\cyru}}1
      {ф}{{\cyrf}}1
      {х}{{\cyrh}}1
      {ц}{{\cyrc}}1
      {ч}{{\cyrch}}1
      {ш}{{\cyrsh}}1
      {щ}{{\cyrshch}}1
      {ъ}{{\cyrhrdsn}}1
      {ы}{{\cyrery}}1
      {ь}{{\cyrsftsn}}1
      {э}{{\cyrerev}}1
      {ю}{{\cyryu}}1
      {я}{{\cyrya}}1
}
\lstset{style=sqlstyle}

\newtheorem{theorem}{Теорема}
\newtheorem{definition}{Определение}

\title{Билет №78 \\ \small Язык SQL. Средства определения и изменения схемы БД, определения ограничений целостности. Контроль доступа. Средства манипулирования данными. (СпБГУ)}
\author{Сериков Владислав Александрович}
\date{16 мая 2026}

\begin{document}

\maketitle
\tableofcontents
\newpage

\section{Введение}
\textbf{SQL (Structured Query Language)} --- декларативный язык программирования, применяемый для создания, модификации и управления данными в реляционных базах данных.

Для обеспечения теоретической базы SQL опирается на реляционную алгебру. Связь между SQL и реляционной алгеброй описывается теоремой о реляционной полноте.

\begin{definition}
Язык запросов называется \textbf{реляционно полным}, если он обладает выразительной мощностью, не уступающей базовой реляционной алгебре (содержащей операции: выборка, проекция, декартово произведение, объединение, разность).
\end{definition}

\begin{theorem}[Реляционная полнота ядра SQL]
Язык SQL, включающий базовый блок \texttt{SELECT-FROM-WHERE}, а также операции \texttt{UNION}, \texttt{EXCEPT} и возможность использования вложенных подзапросов, является реляционно полным.
\end{theorem}

\begin{proof}
Доказательство проведем методом структурной индукции по построению выражения реляционной алгебры. Покажем, что любую из пяти базовых операций можно выразить средствами SQL.

\textbf{База индукции:} Отношение $R$.
В SQL тривиально выражается как: \lstinline|SELECT * FROM R;|.

\textbf{Индукционный переход:} Пусть выражения $E_1$ и $E_2$ уже могут быть выражены SQL-запросами $Q_1$ и $Q_2$. Покажем выразимость базовых операций:
\begin{enumerate}
    \item \textit{Проекция} $\pi_{A_1, \dots, A_n}(E_1)$.
    Транслируется в SQL путем указания нужных атрибутов в секции \texttt{SELECT}:
    \begin{lstlisting}
SELECT DISTINCT A_1, ..., A_n FROM (Q_1) AS T;
    \end{lstlisting}
    (Ключевое слово \texttt{DISTINCT} необходимо для строгого соответствия алгебре, так как она оперирует множествами, а не мультимножествами).

    \item \textit{Выборка (селекция)} $\sigma_P(E_1)$, где $P$ --- предикат.
    Транслируется с помощью секции \texttt{WHERE}:
    \begin{lstlisting}
SELECT * FROM (Q_1) AS T WHERE P;
    \end{lstlisting}

    \item \textit{Декартово произведение} $E_1 \times E_2$.
    \begin{lstlisting}
SELECT * FROM (Q_1) AS T1 CROSS JOIN (Q2) AS T2;
    \end{lstlisting}

    \item \textit{Объединение} $E_1 \cup E_2$.
    \begin{lstlisting}
(Q_1) UNION (Q_2);
    \end{lstlisting}
    (В SQL \texttt{UNION} по умолчанию удаляет дубликаты, что соответствует теоретико-множественному объединению).

    \item \textit{Разность} $E_1 \setminus E_2$.
    \begin{lstlisting}
(Q_1) EXCEPT (Q_2);
    \end{lstlisting}
    (В некоторых СУБД используется оператор \texttt{MINUS}).
\end{enumerate}
Поскольку все базовые операции реляционной алгебры выразимы в SQL, язык является реляционно полным. Ч.т.д.
\end{proof}

\section{Средства определения и изменения схемы БД (DDL)}
\textbf{DDL (Data Definition Language)} --- подмножество SQL, предназначенное для описания структур базы данных.

Основные операторы DDL:
\begin{itemize}
    \item \texttt{CREATE} --- создание объектов БД (таблиц, индексов, представлений).
    \item \texttt{ALTER} --- модификация существующих объектов (добавление/удаление столбцов, изменение типов).
    \item \texttt{DROP} --- удаление объектов из БД.
    \item \texttt{TRUNCATE} --- быстрое удаление всех данных из таблицы с сохранением её структуры.
\end{itemize}

\textbf{Пример создания и изменения таблицы:}
\begin{lstlisting}
CREATE TABLE Students (
    StudentID INT PRIMARY KEY,
    Name VARCHAR(100) NOT NULL,
    EnrollmentYear INT
);

ALTER TABLE Students ADD Email VARCHAR(255);
DROP TABLE Students;
\end{lstlisting}

\section{Ограничения целостности (Integrity Constraints)}
Ограничения целостности гарантируют корректность и непротиворечивость данных. Они определяются в DDL при создании или изменении таблицы.

\begin{enumerate}
    \item \textbf{PRIMARY KEY} (Первичный ключ) --- уникально идентифицирует каждую запись в таблице. Не допускает значений \texttt{NULL}.
    \item \textbf{FOREIGN KEY} (Внешний ключ) --- обеспечивает ссылочную целостность между таблицами. Гарантирует, что значение в дочерней таблице должно присутствовать в родительской. Может сопровождаться каскадными операциями: \texttt{ON DELETE CASCADE}, \texttt{ON UPDATE SET NULL} и др.
    \item \textbf{UNIQUE} --- гарантирует уникальность значений в столбце (допускает одно или несколько значений \texttt{NULL} в зависимости от СУБД).
    \item \textbf{NOT NULL} --- запрещает наличие пустых значений (\texttt{NULL}) в столбце.
    \item \textbf{CHECK} --- проверяет выполнение заданного логического условия для каждой строки.
\end{enumerate}

\textbf{Пример использования ограничений:}
\begin{lstlisting}
CREATE TABLE Enrollments (
    EnrollmentID INT PRIMARY KEY,
    StudentID INT,
    CourseID INT,
    Grade DECIMAL(3,2) CHECK (Grade >= 0.0 AND Grade <= 5.0),
    CONSTRAINT fk_student FOREIGN KEY (StudentID)
        REFERENCES Students(StudentID) ON DELETE CASCADE
);
\end{lstlisting}

\section{Контроль доступа (DCL)}
\textbf{DCL (Data Control Language)} --- подмножество SQL, управляющее правами доступа пользователей к объектам БД.

\begin{itemize}
    \item \texttt{GRANT} --- предоставляет привилегии.
    \item \texttt{REVOKE} --- отзывает привилегии.
\end{itemize}

Синтаксис:
\begin{lstlisting}
GRANT SELECT, INSERT ON Students TO 'user1';
REVOKE DELETE ON Students FROM 'user1';
\end{lstlisting}

Также в современных СУБД используется ролевая модель доступа (RBAC), позволяющая выдавать привилегии не конкретному пользователю, а роли:
\begin{lstlisting}
CREATE ROLE data_entry;
GRANT INSERT, UPDATE ON Enrollments TO data_entry;
GRANT data_entry TO 'alice';
\end{lstlisting}

\section{Средства манипулирования данными (DML)}
\textbf{DML (Data Manipulation Language)} --- подмножество SQL для работы с данными внутри таблиц. Включает \texttt{INSERT}, \texttt{UPDATE}, \texttt{DELETE} и \texttt{SELECT} (часто \texttt{SELECT} выделяют в DQL --- Data Query Language).

\subsection{Модификация данных}
\begin{lstlisting}
-- Добавление данных
INSERT INTO Students (StudentID, Name) VALUES (1, 'Ivan Ivanov');

-- Обновление данных
UPDATE Students SET EnrollmentYear = 2023 WHERE StudentID = 1;

-- Удаление данных
DELETE FROM Students WHERE EnrollmentYear < 2020;
\end{lstlisting}

\subsection{Алгоритм логической обработки DML-запроса (SELECT)}
Для понимания семантики оператора \texttt{SELECT} используется алгоритм логического (концептуального) выполнения запроса. Порядок написания ключевых слов в SQL не совпадает с порядком их логической обработки СУБД.

\textbf{Шаги алгоритма логической обработки запроса:}

\begin{enumerate}
    \item \textbf{FROM (и JOIN):} Определение источников данных. Вычисляется декартово произведение всех указанных таблиц. Если указаны соединения (\texttt{JOIN}), применяются условия объединения (\texttt{ON}), фильтруя декартово произведение.
    \item \textbf{WHERE:} Применение предиката фильтрации к строкам, полученным на шаге 1. Строки, для которых предикат дает \texttt{FALSE} или \texttt{UNKNOWN}, отбрасываются.
    \item \textbf{GROUP BY:} Группировка строк по уникальным значениям указанных столбцов. Каждая группа сворачивается в одну строку.
    \item \textbf{HAVING:} Фильтрация сгруппированных данных (вычисление агрегатных функций и проверка условий для групп).
    \item \textbf{SELECT:} Проекция. Вычисляются выражения, указанные в списке выборки. Назначаются псевдонимы (\texttt{AS}).
    \item \textbf{DISTINCT:} Удаление дубликатов из результирующего набора строк.
    \item \textbf{ORDER BY:} Сортировка итогового набора данных по заданным критериям (направление \texttt{ASC} или \texttt{DESC}).
    \item \textbf{LIMIT / OFFSET:} Пропуск заданного числа строк и ограничение размера вывода (в диалектах SQL Server используется \texttt{TOP}, в Oracle — \texttt{FETCH FIRST}).
\end{enumerate}

\textbf{Иллюстрация и минимальный пример алгоритма:}

Рассмотрим запрос, который выводит номера курсов, на которых обучается более 10 студентов, отсортированные по убыванию числа студентов:

\begin{lstlisting}
SELECT CourseID, COUNT(StudentID) as StudentCount
FROM Enrollments
WHERE Grade >= 3.0
GROUP BY CourseID
HAVING COUNT(StudentID) > 10
ORDER BY StudentCount DESC;
\end{lstlisting}

\textit{Пошаговая иллюстрация:}
\begin{itemize}
    \item \textbf{Шаг 1 (FROM):} Берется полная таблица \texttt{Enrollments}.
    \item \textbf{Шаг 2 (WHERE):} Отбрасываются строки, где оценка меньше $3.0$.
    \item \textbf{Шаг 3 (GROUP BY):} Оставшиеся строки разбиваются на "корзины" (группы) по значению \texttt{CourseID}.
    \item \textbf{Шаг 4 (HAVING):} Для каждой корзины вычисляется \texttt{COUNT(StudentID)}. Корзины, в которых $\leq 10$ элементов, удаляются из рассмотрения.
    \item \textbf{Шаг 5 (SELECT):} Для оставшихся корзин формируются результирующие строки: номер курса и вычисленное количество (создается алиас \texttt{StudentCount}).
    \item \textbf{Шаг 6 (ORDER BY):} Сформированные строки сортируются по убыванию значения \texttt{StudentCount}.
\end{itemize}

\end{document}
